% ID: SEL-D0075-OA-E02
% Tanda S075.
\textbf{a)} En el tramo racional se excluye $x=1$, donde el denominador se anula. El tramo polinómico admite todos los $x\ge2$. Así, $D_f=\boxed{\mathbb R\setminus\{1\}}$.

\textbf{b)} El límite por la izquierda en $2$ es $2\cdot2/(2-1)=4$, mientras que $f(2)=2\cdot2^2-10\cdot2=-12$, que coincide con el límite por la derecha. Como la derivabilidad implica continuidad, $f$ no es derivable en $2$.

\textbf{c)} Cerca de $0$ se usa el tramo racional. Por la regla del cociente,
\[
 f'(x)=\frac{2(x-1)-2x}{(x-1)^2}=-\frac2{(x-1)^2}.
\]
La gráfica pasa por $(0,0)$ y la pendiente allí es $-2$. La tangente es $\boxed{y=-2x}$.
